%page 0049
\chapter{Amphetamine}
\section{From the birth of an immortal substance}

The chemical structure of amphetamine (1) (alphamethylphenethylamine) first appeared in 1887. The Romanian chemist Laxär Edeleanu, active at the Berlin University, described the synthesis for the first time in the article "About some derivatives of phenyl methacrylic acid and phenylisobutyric acid", naming the substance as phenisopropylamine [1]. Interestingly, around four decades then passed before the physiological properties were first reported in the literature. Presumably, the interesting pharmacological properties of amphetamine (1) remained for Edeleanu undiscovered - or he never made it public.

\includegraphics{placeholder}

It is thanks to the Australian historian of science Nicolas Rasmussen that light was shed on the unknown historical beginnings of amphetamine. Before he published his report in 2006 with the title Making the first anti-depressant: Amphetamine in Amenican Medicine, 1929-1950, some inconsistencies were encountered in the explanation of the exact historical course of the development of amphetamine [2]. In addition, time tables from textbooks and lexicons contradicted each other. Even with the use of primary literature, extracted from the SciFinder® database, in addition to the structural search, the names of the researchers working at the time (Piness, Miller, Alles, Hartung, Munch, Prinzmetal) as well as the names of Smith, Kline and French Laboratories (SKF) were entered, it was not possible to reconstruct the exact course. Rasmussen worked his way through various archives and was able to gain access to the SKF archives, where he meticulously shaped documents, letters and memos into a complete picture. The following description of the history of amphetamine is therefore based strongly on Rasmussen's description [2]. It was also supplemented with publications by other researchers.
\bigbreak

At the beginning of the 20th century, adrenaline (2) was considered the most effective remedy for asthmatic complaints, as it dilates the bronchi. One knew how to use this therapeutically in the case of an asthmatic attack. It was possible as early as 1905, relatively pure

\clearpage

%page 0050
\clearpage

%page 0051
\clearpage

%page 0052
\clearpage

%page 0053
\clearpage

%page 0054
\clearpage

%page 0055
\clearpage

%page 0056
\clearpage

%page 0057
\clearpage

%page 0058
\clearpage

%page 0059
\clearpage

%page 0060
\clearpage

%page 0061
\clearpage

%page 0062
\clearpage

%page 0063
\clearpage

%page 0064
\clearpage

%page 0065
\clearpage

%page 0066
\clearpage


